O programa \emph{Artemis}, liderado pela \gls{nasa}, visa estabelecer uma presença de longo prazo no entorno da Lua e em sua superfície \cite{nasaArtemis}.
Esse objetivo demanda capacidades de \gls{rvdb} compatíveis com cenários cislunares, nos quais a maior distância em relação à Terra introduz atrasos de comunicação na ordem de segundos, limitando intervenções imediatas do segmento terrestre \cite{keeports2006}.
Nesse contexto, torna-se necessário ampliar o grau de autonomia dos sistemas responsáveis por aproximação e acoplamento entre veículos e módulos.

Para superar essas restrições, arquiteturas de \gls{rvdb} associadas ao ambiente lunar tendem a incorporar funções de decisão e execução a bordo.
Sensores como LiDAR \cite{farkasvolgyi2023} e câmeras estereoscópicas, aliados a algoritmos de percepção e processamento em tempo real, permitem detectar alvos e estimar estados relativos, apoiando o cálculo de trajetórias de aproximação com menor dependência de supervisão contínua em solo \cite{ibmComputerVision}.

As tecnologias associadas ao \emph{Artemis} também são pertinentes a missões de exploração mais distantes, nas quais atrasos de comunicação e durações de campanha ampliadas impõem requisitos adicionais de autonomia.
Assim, a integração de sistemas autônomos de \gls{rvdb} contribui para operações em ambientes espaciais com restrições de comunicação, reduzindo a necessidade de intervenções diretas e recorrentes por parte de operadores humanos.
